\documentclass[12pt]{article}
\usepackage[utf8]{inputenc}
\usepackage[T1]{fontenc}
\usepackage{hyperref}
\title{Yawuru Language Resources for K-12 Curriculum Development}
\author{Compiled by Claude Sonnet 5.5 (claude-sonnet-5-5) for Andr\'as Kornai}
\date{}
\begin{document}
\maketitle

\section{Introduction}
Yawuru (ISO 639-3 ywr, Glottocode yawu1244) is an Eastern Nyulnyulan language of the Broome area in the Kimberley, Western Australia. It is dormant, with younger people learning it as a second language. This report lists only resources that appeared in search results. No standalone online dictionary was found.

\section{Primary Educational Resources}
The search found no school curriculum. It found a phrasebook, \emph{Yawuru Ngan-ga: a Phrasebook of the Yawuru Language} (Magabala Books, 1995), which Wikipedia cites. I did not find a link to the book itself. The Wikipedia article on the language is here: \url{https://en.wikipedia.org/wiki/Yawuru_language}

\section{Linguistic and Reference Materials}
\begin{itemize}
\item AIATSIS Austlang record K1 for Yawuru, which includes Hosokawa's 1991 thesis and his 2011 Lincom Europa grammar among its sources: \url{https://collection.aiatsis.gov.au/austlang/language/K1}
\item AIATSIS collection page for Yawuru: \url{https://collection.aiatsis.gov.au/node/1508637}
\item Omniglot overview of Yawuru: \url{https://omniglot.com/writing/yawuru.htm}
\end{itemize}

\section{Religious and Community Resources}
The search turned up the Walalangga Yawuru Ngan-ga language program, which runs adult classes, and Nyamba Buru Yawuru, which Wikipedia cites as a source of language information. I did not see a URL for either in the results, so none is given here. The community's own article is at \url{https://en.wikipedia.org/wiki/Yawuru}.

\section{Digital and Audio Resources}
The Endangered Languages Project has a sample of Yawuru vocabulary about the environment: \url{https://endangeredlanguages.com/lang/3160/samples/9821?hl=en}

\section{Curriculum Development Context}
No evidence was found that Yawuru is taught in schools. Any revitalization curriculum should be led by Yawuru people and their language program. The published grammar, the phrasebook and the AIATSIS holdings are the main sources the search found.
\end{document}